(1) 在共转系下单分子有效离心势 \(\Phi_{\rm eff}=-m\Omega^2r^2/2\)。等温理想气体满足 \(dp/dr=mn(r)\Omega^2r\)，而 \(p=nk_BT\)，故 \(d\ln n/dr=m\Omega^2r/(k_BT)=2\alpha r\)。积分得到
\[
\boxed{n(r)=n_0e^{\alpha r^2}},
\]
密度向外增加，是离心稳定平衡，而非把离心力错写向内得到指数衰减。必须有有限半径壁面，否则指数分布在无限平面不可归一。

(2) 总粒子数
\(N=2\pi Hn_0\int_0^R re^{\alpha r^2}dr=\pi Hn_0(e^{\alpha R^2}-1)/\alpha\)，所以
\[
n_0=\frac{N\alpha}{\pi H(e^{\alpha R^2}-1)},\quad
p(R)=k_BTn_0e^{\alpha R^2}.
\]
当 \(\alpha\to0\) 时利用 \(e^{\alpha R^2}-1\simeq\alpha R^2\) 恢复 \(n_0=N/(\pi R^2H)\)，壁压为 \(Nk_BT/(\pi R^2H)\)；不可将 \(\alpha=0\) 直接代入未化简的零比零式。

(3) 用截面积权重 \(2\pi r\,dr\) 计算
\[
\langle r^2\rangle
=\frac{\int_0^Rr^3e^{\alpha r^2}dr}{\int_0^Rr e^{\alpha r^2}dr}
=\boxed{\frac{R^2e^{\alpha R^2}}{e^{\alpha R^2}-1}-\frac1\alpha}.
\]
气体平均绕轴角动量 \(L_z=m\Omega N\langle r^2\rangle\)。小 \(\Omega\) 时 \(\langle r^2\rangle\to R^2/2\)，故 \(L_z\simeq\frac12NmR^2\Omega\)；强旋转 \(\alpha R^2\gg1\) 时 \(\langle r^2\rangle\simeq R^2-1/\alpha\)，气体集中在近壁薄壳而不是中心。

(4) 在旋转热浴的平衡配分函数中，单粒子Boltzmann权为 \(\exp[-\beta(H_{\rm lab}-\Omega L_z)]\)。对动量完成平方，积分因子给 \(\exp[\beta m\Omega^2r^2/2]=e^{\alpha r^2}\)，正好再现第一问。在固定 \(T,N,R,H\) 下 \(F_{\rm rot}=-k_BT\ln Z_{\rm rot}\) 且
\(-\partial_\Omega F_{\rm rot}=k_BT\partial_\Omega\ln Z_{\rm rot}=\langle L_z\rangle\)。
因此共转势是一种Legendre势而非实验室的额外真实能量来源；容器施加约束并维持转速时机械功应在实验室账本中另算。
